\begin{pastexam}[label={px:ncu-110-dsa-07}]{中央資工所 110 學年度 複選題第 7 題}
（複選；每題 5 分，全對才給分，答錯一個選項扣 1.5 分）

Insert the set \{23, 12, 34, 46, 28, 11, 6, 7, 0, 33, 30, 45\} into an empty hash Table T of size
17. Suppose that the hash function \textbf{h}(k , i) = (k + i) mod 17, where i = 0 to 16. If we
remove 46 after inserting the set, which keys should modify their positions?\\
(A)~11 \quad (B)~0 \quad (C)~33 \quad (D)~45 \quad (E)~30
\end{pastexam}

\begin{pastsoln}
答案：(A)(D)(E)。

$h(k, i) = (k + i) \bmod 17$ 是線性探測（資結書第~12.3~節）。 題目問刪除 $46$ 後哪些鍵要
「換位置」，所以這裡的刪除不是放刪除標記（資結書第~12.3.4~節；放標記的話沒有鍵會移動），而是
真的把格子清空，再把同一個叢集裡、排在它後面的鍵重新放好，讓之後的搜尋不會被空格切斷。
集合依寫出的順序插入。

\textbf{插入}。 每個鍵從 $k \bmod 17$ 開始往後找空格：
\begin{center}
\small
\begin{tabular}{cccl}
\toprule
鍵 & $k \bmod 17$ & 放在 & 經過的格子 \\
\midrule
$23$ & $6$ & $6$ & \\
$12$ & $12$ & $12$ & \\
$34$ & $0$ & $0$ & \\
$46$ & $12$ & $13$ & $12$ 有 $12$ \\
$28$ & $11$ & $11$ & \\
$11$ & $11$ & $14$ & $11, 12, 13$ 都有鍵 \\
$6$ & $6$ & $7$ & $6$ 有 $23$ \\
$7$ & $7$ & $8$ & $7$ 有 $6$ \\
$0$ & $0$ & $1$ & $0$ 有 $34$ \\
$33$ & $16$ & $16$ & \\
$30$ & $13$ & $15$ & $13, 14$ 有鍵 \\
$45$ & $11$ & $2$ & $11$ 到 $16$、$0$、$1$ 都有鍵，繞回來 \\
\bottomrule
\end{tabular}
\end{center}
插入完的表：
\begin{center}
\small
\setlength{\tabcolsep}{3.5pt}
\begin{tabular}{l*{17}{c}}
\toprule
格 & 0 & 1 & 2 & 3 & 4 & 5 & 6 & 7 & 8 & 9 & 10 & 11 & 12 & 13 & 14 & 15 & 16 \\
\midrule
鍵 & 34 & 0 & 45 & & & & 23 & 6 & 7 & & & 28 & 12 & 46 & 11 & 30 & 33 \\
\bottomrule
\end{tabular}
\end{center}

\textbf{刪除 $46$}。 把第 $13$ 格清空。 從第 $11$ 格開始到第 $2$ 格（過了第 $16$ 格繞回第 $0$ 格）是一整個
連續的叢集，第 $13$ 格後面的鍵依序是 $11, 30, 33, 34, 0, 45$，到第 $3$ 格空格為止。 把它們依序
拿出來、照同一個雜湊函數重新插入：
\begin{itemize}
\item $11$（在第 $14$ 格）：從 $11$ 找，$11$、$12$ 有鍵，第 $13$ 格空了，搬到 $13$。
\item $30$（在第 $15$ 格）：從 $13$ 找，$13$ 剛放了 $11$，第 $14$ 格空了，搬到 $14$。
\item $33$（第 $16$ 格）：本來就在 $33 \bmod 17 = 16$，不動。
\item $34$（第 $0$ 格）、$0$（第 $1$ 格）：從 $0$ 找，位置不變。
\item $45$（在第 $2$ 格）：從 $11$ 找，$11$ 到 $14$ 都有鍵，第 $15$ 格空了，搬到 $15$。
\end{itemize}
結果第 $13, 14, 15$ 格是 $11, 30, 45$，第 $2$ 格空出來。 換位置的是 $11$、$30$、$45$，選 (A)(D)(E)；
(B) $0$ 和 (C) $33$ 都在原位。

用另一種常見的做法「往前補洞」也得到同樣的結果：從空格往後掃，遇到一個鍵，如果它的起始位置
不在「空格之後、它目前位置之前（含）」這一段裡，就把它搬進空格，空格移到它原來的位置。 $11$
（起始 $11$）、$30$（起始 $13$）、$45$（起始 $11$）會被搬，$33$、$34$、$0$ 的起始位置就在那一段裡，
不搬。
\end{pastsoln}
